% !TeX root = ../main.tex
\section{Reproducibility}
\label{app:repro}

This appendix lists artifacts, evaluation commands, and environment details for
reproducing the numbers cited in Section~\ref{sec:results}. A machine-readable
checklist is maintained in \texttt{phase5/REPRODUCIBILITY.md}.

\subsection{Dataset and splits}

We use CESNET-QUIC22-XS~\cite{cesnet_quic22}: QUIC flows collected on a national
research-and-education backbone. Training uses calendar week W-2022-44; testing uses
the subsequent week W-2022-45 (49{,}305 flows, 64 application classes after filtering
classes with fewer than 200 training samples). Each flow is a $(3,30)$ tensor of IPT,
direction, and normalized packet size.

\subsection{Checkpoints and artifacts}

Table~\ref{tab:repro} summarizes primary files. Obfuscated tensors are stored as
\path{phase3/artifacts/obfuscated_<policy>.pt}; overhead is logged in
\path{phase3/artifacts/obfuscation_manifest.json}.

\begin{table}[t]
  \centering
  \caption{Key reproducibility artifacts.}
  \label{tab:repro}
  \tiny
  \setlength{\tabcolsep}{2pt}
  \renewcommand{\arraystretch}{1.12}
  \begin{tabular}{@{}>{\raggedright\arraybackslash}p{0.22\columnwidth}>{\raggedright\arraybackslash}p{0.74\columnwidth}@{}}
    \toprule
    \textbf{Item} & \textbf{Relative path} \\
    \midrule
    Test tensors & \texttt{phase1/artifacts/test\_tensors.pt} \\
    Transf.\ ckpt & \texttt{phase2/artifacts/transformer\_production.pt} \\
    BiLSTM ckpt & \texttt{phase2/artifacts/cnn\_bilstm\_best.pt} \\
    Obfuscated & \texttt{phase3/artifacts/obfuscated\_*.pt} \\
    Results CSV & \texttt{phase4/results/accuracy\_results.csv} \\
    Bootstrap CIs & \texttt{phase4/bootstrap\_ci.csv} \\
    \bottomrule
  \end{tabular}
\end{table}

\subsection{Evaluation commands}

\textbf{Tier A/B (Transformer sweep):}
\begin{center}
\fbox{%
  \begin{minipage}{0.94\columnwidth}
  \scriptsize\ttfamily\raggedright
  cd phase4\par
  python run\_experiments.py\par
  \hspace{1em}--device cuda\par
  \hspace{1em}--checkpoint\par
  \hspace{2em}../phase2/artifacts/\par
  \hspace{2em}transformer\_production.pt
  \end{minipage}}
\end{center}

\textbf{Tier C (bootstrap / ablation / arch.):}
\begin{center}
\fbox{%
  \begin{minipage}{0.94\columnwidth}
  \scriptsize\ttfamily\raggedright
  python run\_tier\_c.py --all\par
  \hspace{1em}--device cuda\par
  \hspace{1em}--checkpoint\par
  \hspace{2em}../phase2/artifacts/\par
  \hspace{2em}transformer\_production.pt\par
  python run\_tier\_c.py --bilstm\par
  \hspace{1em}--compare-architectures\par
  \hspace{1em}--device cuda\par
  \hspace{1em}--bilstm-checkpoint\par
  \hspace{2em}../phase2/artifacts/\par
  \hspace{2em}cnn\_bilstm\_best.pt
  \end{minipage}}
\end{center}

Re-running \texttt{run\_tier\_c.py --all} regenerates bootstrap CIs, paired tests, and
channel-ablation tables referenced in Section~\ref{sec:results}.

\subsection{Environment and ethics}

Experiments were run with pinned dependencies in \texttt{phase4/requirements.txt}.
CESNET-QUIC22 is distributed under the Zenodo terms cited in~\cite{cesnet_quic22};
we use only flow metadata (sizes, directions, timing) and application labels---no
payload content is accessed or stored in our pipeline.
